\chapter{1949:William F. Giauque}

\label{chap:chemistry1949}

The Nobel Prize in Chemistry for the year 1949 was awarded to William F. Giauque.

\textbf{Motivation:}

for his contributions in the field of chemical thermodynamics, particularly concerning the behaviour of substances at extremely low temperatures

\section{William F. Giauque}

\subsection*{Early Life and Education}

William F. Giauque was born on 1895-05-12 in Niagara Falls, Ontario, Canada.

\subsection*{Career and Major Research}

Giauque moved to the United States as a young man and studied at the University of California, Berkeley, where he earned his doctorate in 1922 and spent his entire scientific career. With Herrick Johnston he discovered the isotopes oxygen-17 and oxygen-18 in 1929, work that showed oxygen could no longer serve as the sole basis of the atomic weight scale.

\subsection*{Nobel Prize-winning Work}

The work for which William F. Giauque was awarded the Nobel Prize in Chemistry in 1949 was described by the Nobel Committee as: "for his contributions in the field of chemical thermodynamics, particularly concerning the behaviour of substances at extremely low temperatures".

Working with D. P. MacDougall, Giauque applied adiabatic demagnetization of a paramagnetic salt to reach temperatures below one degree absolute, attaining about 0.25 K in 1933. The method, proposed independently by Giauque and Peter Debye, made possible the study of entropy and heat capacity of substances at the lowest attainable temperatures.

\subsection*{Significance and Impact}

The demagnetization method and Giauque's low-temperature calorimetry provided critical tests of the third law of thermodynamics. His careful measurements of entropies helped resolve discrepancies between statistical and thermal entropies of gases.

\subsection*{Later Career and Honors}

Giauque remained at Berkeley, where he continued research in low-temperature thermodynamics. He died in 1982.

\subsection*{Legacy}

Adiabatic demagnetization, introduced by Giauque, remains the standard means of reaching temperatures in the millikelvin range, and his thermodynamic methods shaped modern low-temperature physics and chemistry.

\subsection*{Selected Publications}

\begin{itemize}

\item \emph{Attainment of Temperatures Below 1\textdegree{} Absolute by Demagnetization of Gd2(SO4)3\textperiodcentered 8H2O}, Physical Review, 1933.

\end{itemize}

\clearpage

\section*{Chapter Summary}

The 1949 prize recognized contributions to chemical thermodynamics, especially the behaviour of substances at extremely low temperatures. William F. Giauque's adiabatic demagnetization technique reached temperatures near absolute zero and enabled precise tests of the third law of thermodynamics.
